\begin{enumerate}[label=\textbf{\ansnum.\arabic*}]
  \item First case: the best is robot $1$ itself ($1.4<1.5$), so it stays striker. Second case: robot $2$ is nearer ($1.05$), but $1.05<1.4-0.3=1.1$ is true, so the striker becomes robot $2$ (the margin of $0.35$\,m is larger than $0.3$). With the value $1.15$ instead the striker would stay robot $1$.
  \item $12000/(3\cdot1200)=3.33$\,Hz per robot. With $5$ robots at $1$\,Hz the budget lasts $12000/5=2400$\,s, so there are $2400-1200=1200$\,s of \emph{ready}/\emph{set} time to spare (which is much more than needed).
  \item $12000/5=2400$\,s. After $600$\,s: $12000-5\cdot600=9000$ messages are left.
  \item $6000/(5\cdot2)=600$\,s, i.e. $10$ minutes, twice the length of the $5$ minutes of extra time of the draft rules. In extra time the budget per second is larger than in the main game, so there is no need to economise there.
  \item Example: player number ($1$ byte), pose $x,y$ in centimetres ($2+2$ bytes) and heading in $1/100$ rad or in degrees ($2$), ball $x,y$ ($2+2$), ball age in $0.1$\,s ($1$), role ($1$), fallen flag ($1$, can share a byte), time stamp ($4$): $18$ bytes. With protobuf (field numbers and varints) a similar message takes about $25$--$40$ bytes; the exact size depends on the values, because small numbers use fewer bytes. Both are far below $512$.
  \item Ball holding: a timer that starts when the ball is in the dribbling zone; at $3$\,s the robot must kick or pass. Local game stuck: if the ball has not moved for $6$\,s the striker kicks in the direction of the opponent's goal. Pushing: obstacle avoidance that keeps a distance and a rule that the robot stops pushing and walks around if there is contact for more than $2$--$3$\,s.
  \item The decision must be consistent: if every robot works with its own noisy ball and its own positions, two robots can both think that they are nearest, both go for the ball and block each other (or none goes). With the same data (shared via the messages, or one robot as ``decider'') and the same rule the result is the same. The hysteresis state (who is the current striker) is also shared.
  \item $3\times6$\,m. The area is $3$\,m deep: it ends at $3$\,m from the goal line, and is $6$\,m wide ($3$\,m to each side of the centre line of the goal). The goal area ($1\times4$\,m) is inside it.
  \item With equal true distances and symmetric noise the two robots are equally likely to be nearest. In two independent measurements the nearest robot changes with probability $1/2$ ($1/2\cdot1/2\cdot2$). With hysteresis $h=0.3$\,m and noise $0.15$\,m the difference of two measurements has standard deviation $0.15\sqrt2=0.21$\,m, so a switch needs the difference to be larger than $0.3$\,m: probability $P(Z>0.3/0.21)=P(Z>1.41)=0.079$ per decision in that direction. (This is lower than $0.5$ but still not small at $10$ decisions per second; the real simulation shows $13.6$ changes in $60$\,s.)
  \item A Python simulation as in the chapter (the program for the figure was run: the numbers for the mean number of role changes are $115.7$, $60.8$, $13.6$, $6.6$ for $h=0,0.1,0.3,0.5$ and $6$ for the noise-free decision). The delay grows with $h$: it is the time during which the new nearest robot is better by less than $h$. (I did not measure the delay; this is the task of the problem.)
\end{enumerate}
